% !TEX program = lualatex
% Transparencias de Fundamentos de las Matemáticas II
% Clase 13: Números reales I

\documentclass[aspectratio=169,11pt]{beamer}

\usepackage{fontspec}
\usepackage[spanish,es-nodecimaldot]{babel}
\usepackage{amsmath,amssymb,mathtools}
\usepackage{booktabs}
\usepackage{csquotes}
\usepackage{xcolor}

\usetheme{Madrid}
\usecolortheme{default}
\setbeamertemplate{navigation symbols}{}
\setbeamertemplate{footline}[frame number]

\definecolor{FdMBlue}{RGB}{20,55,100}
\definecolor{FdMRed}{RGB}{130,30,30}
\definecolor{FdMGreen}{RGB}{25,100,60}

\setbeamercolor{structure}{fg=FdMBlue}
\setbeamercolor{title}{fg=white,bg=FdMBlue}
\setbeamercolor{frametitle}{fg=white,bg=FdMBlue}
\setbeamercolor{block title}{fg=white,bg=FdMBlue}
\setbeamercolor{block body}{bg=blue!3}
\setbeamercolor{alerted text}{fg=FdMRed}

\setmainfont{Latin Modern Roman}
\setsansfont{Latin Modern Sans}
\setmonofont{Latin Modern Mono}

\newcommand{\N}{\mathbb{N}}
\newcommand{\Z}{\mathbb{Z}}
\newcommand{\Q}{\mathbb{Q}}
\newcommand{\R}{\mathbb{R}}
\newcommand{\C}{\mathbb{C}}
\newcommand{\Pow}{\mathcal{P}}
\newcommand{\id}{\operatorname{id}}
\newcommand{\mcd}{\operatorname{mcd}}
\newcommand{\mcm}{\operatorname{mcm}}
\newcommand{\im}{\operatorname{Im}}

\title[FdM II -- Clase 13]{Números reales I}
\subtitle{Irracionales, completitud e intervalos}
\author{Antonio Falcó Montesinos}
\institute{Universidad CEU Cardenal Herrera}
\date{Fundamentos de las Matemáticas II}

\begin{document}

\begin{frame}
  \titlepage
\end{frame}

\begin{frame}{Objetivos de la clase}
\begin{itemize}
  \item Explicar por qué \(\mathbb{Q}\) es insuficiente.
  \item Introducir \(\mathbb{R}\) como sistema completo.
  \item Trabajar con intervalos.
\end{itemize}
\end{frame}

\begin{frame}{Mapa de la clase}
\tableofcontents
\end{frame}

\section{De \(\mathbb{Q}\) a \(\mathbb{R}\)}

\begin{frame}{De \(\mathbb{Q}\) a \(\mathbb{R}\)}
\begin{block}{Insuficiencia}
\begin{itemize}
  \item La diagonal del cuadrado unidad tiene longitud \(\sqrt{2}\).
  \item \(\sqrt{2}\notin\mathbb{Q}\).
  \item Los racionales son densos, pero tienen huecos.
\end{itemize}
\end{block}
\end{frame}

\begin{frame}{De \(\mathbb{Q}\) a \(\mathbb{R}\)}
\begin{block}{Irracionales}
\begin{itemize}
  \item Un número real \(x\) es irracional si \(x\notin\mathbb{Q}\).
  \item Ejemplos: \(\sqrt{2}\), \(\sqrt{3}\), \(\pi\).
  \item No se escriben como cociente de enteros.
\end{itemize}
\end{block}
\end{frame}

\begin{frame}{De \(\mathbb{Q}\) a \(\mathbb{R}\)}
\begin{block}{Completitud}
\begin{itemize}
  \item Trabajaremos con \(\mathbb{R}\) como cuerpo ordenado completo.
  \item La completitud expresa que los subconjuntos no vacíos y acotados superiormente tienen supremo.
  \item Esta propiedad distingue \(\mathbb{R}\) de \(\mathbb{Q}\).
\end{itemize}
\end{block}
\end{frame}

\section{Intervalos}

\begin{frame}{Intervalos}
\begin{block}{Intervalos acotados}
\begin{itemize}
  \item \((a,b)=\{x\in\mathbb{R}:a<x<b\}\).
  \item \([a,b]=\{x\in\mathbb{R}:a\leq x\leq b\}\).
  \item Los extremos incluidos se indican con corchetes.
\end{itemize}
\end{block}
\end{frame}

\begin{frame}{Intervalos}
\begin{block}{Intervalos no acotados}
\begin{itemize}
  \item \((a,+\infty)\) representa los reales mayores que \(a\).
  \item \((-infty,a]\) representa los reales menores o iguales que \(a\).
  \item \(+\infty\) y \(-\infty\) no son números reales.
\end{itemize}
\end{block}
\end{frame}

\section{Profundización}

\begin{frame}{Profundización}
\begin{block}{Densidad frente a completitud}
\begin{itemize}
  \item Densidad: entre dos racionales hay otro racional.
  \item Completitud: no hay huecos respecto al supremo.
  \item \(\mathbb{Q}\) es denso, pero no completo.
\end{itemize}
\end{block}
\end{frame}

\begin{frame}{Profundización}
\begin{block}{Ejemplo}
\begin{itemize}
  \item El conjunto \(A=\{q\in\mathbb{Q}:q^2<2, q>0\}\) está acotado en \(\mathbb{Q}\).
  \item Su supremo natural es \(\sqrt{2}\).
  \item Pero \(\sqrt{2}\notin\mathbb{Q}\).
\end{itemize}
\end{block}
\end{frame}

\begin{frame}{Profundización}
\begin{block}{Control}
\begin{itemize}
  \item ¿Es \((1,3]\) un subconjunto de \(\mathbb{R}\)?
  \item ¿Contiene a \(1\)?
  \item ¿Contiene a \(3\)?
\end{itemize}
\end{block}
\end{frame}


\section{Detalle y práctica}

\begin{frame}{Detalle y práctica}
\begin{block}{Huecos de \(\mathbb{Q}\)}
\begin{itemize}
  \item Los racionales son densos, pero no completan la recta.
  \item La ecuación \(x^2=2\) no tiene solución en \(\mathbb{Q}\).
  \item Los reales completan esos huecos mediante una noción rigurosa de límite o de corte.
\end{itemize}
\end{block}
\end{frame}

\begin{frame}{Detalle y práctica}
\begin{block}{Actividad breve}
\begin{itemize}
  \item Recordar la prueba de que \(\sqrt{2}\notin\mathbb{Q}\).
  \item Dar dos racionales entre \(1.41\) y \(1.42\).
  \item Explicar por qué densidad no es lo mismo que completitud.
\end{itemize}
\end{block}
\end{frame}

\begin{frame}{Cierre}
\begin{itemize}
  \item Los reales completan los racionales.
  \item Los irracionales no son fracciones de enteros.
  \item Los intervalos formalizan regiones de la recta real.
\end{itemize}
\end{frame}

\begin{frame}{Trabajo recomendado}
\begin{itemize}
  \item Releer las secciones correspondientes del manual.
  \item Identificar definiciones, símbolos nuevos y enunciados principales.
  \item Preparar dudas y ejemplos para el seminario asociado.
\end{itemize}
\end{frame}

\end{document}
